\section{提案システム}
\begin{figure}[htbp]
\begin{center}
\includegraphics[width=40mm]{img/3/system.png}
\caption{提案システムのイメージ図.}
\label{system}
\end{center}
\end{figure}
\vspace{-1.5zh}
%提案手法では, アプリケーションの機能を一つ一つきちんと説明する
我々は, プロのストーリークリエーターを対象とした, ストーリー創作支援を行う web アプリケーションを開発した. 本アプリケーションは, 手塚治虫の漫画『ブラック・ジャック』\cite{ref:Tezuka77}における物語構造分析や, キャラクター分析により得られたデータを LLM の制御に活用する. 機能は大きく2つのセクションに分かれ, 一つはプロットの生成を行うセクション, もう一つはプロットの編集を行うセクションである. 分析に基づく『ブラック・ジャック』の新作プロットの生成を, 直感的な操作により実現するインターフェースと, プロットの改善を目的とした LLM とのインタラクション手法を提供する. いずれのセクションにおいても, ストーリークリエーターと LLM との効果的なやりとりを促すフレームワークを実装した.

%インターフェースを通じた直感的な操作により, 分析に基づく『ブラック・ジャック』の新作プロットの生成を可能とし, プロットの質改善を目的とした LLM とのインタラクション手法を提供する.

%分析に基づいた『ブラック・ジャック』の新作プロットの生成を, インターフェースを通じた直感的な操作により実現し, LLM との効果的なインタラクション手法を提供する.

%我々は, プロのストーリークリエーターを対象とした, ストーリー創作支援を行う web アプリケーションを開発した. 本アプリケーションは, 手塚治虫の漫画『ブラック・ジャック』\cite{ref:Tezuka77}における物語構造分析や, キャラクター分析により得られデータを LLM の制御に活用する. ユーザーは本アプリケーションのインターフェースを通じて, 直感的な操作により, 分析に基づく『ブラック・ジャック』の新作プロットを生成するこができる.

%また, 本アプリケーションは, プロットの改善を目的とする LLM とのインタラクティブな共創手法を提供し, クリエーターのストーリー創作を支援する.

%我々は, プロのストーリークリエーターを対象とした, ストーリー創作支援を行う web アプリケーションを開発した. 本アプリケーションは, 手塚治虫の漫画『ブラック・ジャック』\cite{ref:Tezuka77}における物語構造分析や, キャラクター分析により得られデータを LLM の制御に活用し, クリエーターの物語生成過程を容易にする. ユーザーは, 本アプリケーションのインターフェースを通じて, 直感的な操作により, 分析に基づく『ブラック・ジャック』の新作のプロットを生成するこができる.  また, 本アプリケーションは, プロットの改善を目的とする LLM とのインタラクティブな共創手法を提供し, クリエーターのストーリー創作を支援する. 

%直感的な操作を通じて, 分析に基づく『ブラック・ジャック』の新作のプロットを容易に生成することができる. また, 本アプリケーションは, プロットの改善を目的とする LLM とのインタラクティブな共創手法を提供し, クリエーターのストーリー創作を支援する. 

%我々は, プロのストーリークリエーターのストーリー創作を支援する web アプリケーションを開発した. このアプリケーションは, 手塚治虫原作の漫画『ブラック・ジャック』[手塚治虫 77]に対して行われた物語構造分析や, キャラクターの分析データを活用し, クリエーターの LLM を用いたストーリー生成を容易にする. ユーザーは, このアプリケーションを用いることで, 簡単な選択と入力により, 分析に基づいた『ブラック・ジャック』の新作プロットを生成することができる. また, このアプリケーションは, プロットの改善を目的とする LLM とのインタラクティブな共創手法を提供し, クリエーターのストーリー創作を支援する. 

%機能は大きく2つのセクションに分かれ, 1つはプロットの生成を行うセクション, もう一つはプロットの編集を行うセクションである. いずれのセクションにおいても, ストーリークリエーターと LLM との効果的なやりとりを促すフレームワークを提供する. 

本アプリケーションに接続した LLM は GPT-4\cite{ref:OpenAI23} である. GPT-4 は, OpenAI 社が開発した LLM であり, あらゆる自然言語処理タスクにおいて高い性能を発揮する\cite{ref:OpenAI23}. API も公開されていることから, インターフェースへの接続も容易であり, 本実験では GPT-4 を用いることとした.

\subsection{物語構造}
物語構造とは, 物語が進行するパターンを構造化した概念を指す. この研究では, 約300種のロシアの魔法昔話を31の要素の組み合わせにより記述できると主張したProppの研究\cite{ref:Propp64} などがよく知られている. 

%物語構造とは, 物語が進行するパターンを構造化した概念を指し, 物語の構造化は, 物語の特性や構成を対象とする学問分野, 物語論(narratology)の中で取り扱われてきた. この分野における研究では, 約300種のロシアの魔法昔話を31の要素の組み合わせにより記述できると主張したProppの研究\cite{ref:Propp64} などがよく知られている. 

物語の自動生成に関する研究において, 物語構造を利用した研究はこれまでもさまざま行われてきた. Proppが見出した31機能を並べ替え, プロットの自動生成を試みる研究\cite{ref:Jaya11}などである. しかしながら, Proppが見出した31の要素はあくまで特定のジャンル, ロシアの魔法昔話の分析によって得られたものであり, 一般的な物語の作成に利用することはできないと考えられる. 

そこで村井\cite{ref:Murai22}は, 特定ジャンルの物語に限定されない, 一般的な物語の展開パターンを抽出する目的として, 『ブラック・ジャック』を分析対象とし, 物語構造分析を行った. 村井の研究によると『ブラック・ジャック』は, 複数ジャンルの物語パターンを含む短編集であり, 分析可能な話数も十分である. 加えて作品自体の評価も高い. 一般的な物語の展開パターンを抽出するにあたり, 異なる作者, 異なる作品の様々なジャンルの物語を組み合わせて分析することも考えられるが, しかしその場合, 作者や作品の違いによって, 演出上の表現技法などが大きく異なることとなり, 分析が困難になると考えられる. このような理由により, 同一作者, 同一作品内において, 複数ジャンルの物語を含み, 国内外で高い評価を持つ『ブラック・ジャック』は一般的な物語構造の分析に適した作品であるといえる\cite{ref:Murai22}.

そこで本研究では, 村井の物語構造分析により得られた『ブラック・ジャック』の物語構造, 展開パターンを LLM の制御に用いることとし, 『ブラック・ジャック』の新作プロットの創作支援に取り組んだ. 

『ブラック・ジャック』の新作プロットの創作支援により得た知見は, 同作品が複数ジャンルの物語構造を含む複合的な作品であるため, 一般的な物語作品を生成するための汎用的な手法に貢献すると考えられる.

\subsection{プロットの生成}
プロット生成セクションでは, 図\ref{flow}に示すフローに沿って操作を進める.
\begin{figure}[htbp]
\begin{center}
\includegraphics[width=65mm]{img/3/flow.png}
\caption{プロット生成セクションにおけるフロー.}
\label{flow}
\end{center}
\end{figure}

フローに沿って進めることで, 『ブラック・ジャック』の分析に基づくプロット生成を行うことができる. どのステップにおいても操作は簡潔で, 画面の指示に従い選択と入力を行うことにより容易に次のステップへと移行できるようになっている.

ここでいうプロットとは, 物語の全体のおおまかな流れを表す文章のことであり, 本アプリケーションで生成することができるプロットは, 1000字から 3000字程度のプロットである.

物語そのものではなくプロットを生成する理由は二つある. 一つは, 物語そのものには, 構造の上に演出上のカスタマイズがなされるため, 物語構造による制御の効果を直接活かせないと考えられること. もう一つは, 本アプリケーションがクリエーターの代替となるものではなく, クリエーターの創作支援を目的としたものであるため, 創造性を掻き立てる形式としてプロットの形式を選択したことである.
\subsubsection{ジャンル・テーマの設定}
ここでは, 物語のジャンルおよびテーマを, 『ブラック・ジャック』の分析により得た, 「SF」や「人間ドラマ」など, ジャンル64種と, 「親子愛」や「友情」など, テーマに関するキーワード152種の中から複数選び設定することができる.

分析に基づいた選択肢をあらかじめ用意しておくことで, 創作の手間を削減し, 多岐にわたるキーワードの提示により, 創造性を刺激する. 

選択は, 自身の意思により行うこともできれば, ランダム選択機能を用いて行うこともできる. 選び取る時間の削減とランダムネスによる創造性支援を行う\cite{ref:Osone21}.

\subsubsection{象徴表現の生成}
ここでは, 選択したテーマや自身で考えたテーマの象徴表現を生成することができる. 

ここで言う象徴表現とは, 鳩で平和を表すように, 抽象的な概念を別のものに置き換え表す表現技法のことである. この表現技法は, 小説や映画, 詩など様々な場面で用いられ, 物語に何層もの意味を加えて作品を豊かにするのに役立つものである\cite{ref:Norva23}.

象徴表現を生成することで, プロットに深みを与える手段や気づき, ヒントを獲得することができる.

象徴表現の生成例:

\begin{table}[hbtp]
\small
\begin{center}
{
\label{adopt_t}
\begin{tabular}{cc} 
信頼  →& 手を繋ぐ行為
\end{tabular}
}
\end{center}
\end{table}

\vspace{-3zh}

\subsubsection{テーマの言い換えの生成}
ユーザーは, テーマの言い換えの生成を行うことで, 単語形式で設定したテーマを文章の形式で獲得することができる. 
% \begin{table}[hbtp]
% \begin{center}
% {\small
% \caption{言い換えの生成例}
% \label{themachange}
% \begin{tabular}{cc}\hline
% テーマ & 言い換え\\ \hline
% 親子愛  &\begin{tabular}{c}
% 親と子が互いの理解や共感を深め、\\困難を乗り越える過程を描く物語
% \end{tabular} \\ \hline
% \end{tabular}
% }
% \end{center}
% \end{table}
% \begin{center}
% \begin{minipage}[htbp]{70mm}
% \small
% 「親子愛」 → 「親と子が互いの理解や共感を深め、困難を乗り越える過程を描く物語」
% \end{minipage}
% \end{center}

LLM は物語を生成する際, 入力として与えた単語をそのままの形で登場させてしまう傾向がある. しかし, 物語のテーマは, 作品全体を通して伝わることが望ましく, 生成するプロットの中に直接単語として現れてしまうことは好ましくない.

テーマの言い換えの生成により, 文章形式でテーマを設定し直すことで, 選択したキーワードが生成するプロットの中に直接出てくることを防ぐことができる.

テーマの言い換えとして生成された文章は, ユーザー自身の手で書き換えることもでき, LLM とクリエーターとの共創を実現する.

言い換えの生成例:
\begin{table}[hbtp]
\begin{center}
{\footnotesize
\label{adopt_t}
\begin{tabular}{cc} 
親子愛  →&\begin{tabular}{c}
親と子が互いの理解や共感を深め、\\困難を乗り越える過程を描く物語\end{tabular} 
\end{tabular}
}
\end{center}
\end{table}
\vspace{-1.5zh}

なお, 同じキーワードに対する言い換えの文章が同じ文章になるとは限らない. 

\subsubsection{キャラクターの設定}
ここでは, 物語に登場させる, 『ブラック・ジャック』の既存キャラクターと, 自身で作り上げる新規キャラクターの設定を行うことができる.

既存キャラクターには, 「ブラック・ジャック」や「ピノコ」など, 漫画『ブラック・ジャック』に登場する主要なキャラクターを選択肢として用意しており, ユーザーはその中からプロットに登場させるキャラクターを選ぶことができる. それぞれのキャラクターに対し, 「主人公」や「協力者」など物語上での役割や性格といった属性を, 分析に基づいた選択肢の中から設定することができ, 既存キャラクターの言動が原作からかけ離れたものにならぬよう工夫を施している.

\begin{figure}[htbp]
\begin{center}
\includegraphics[width=81mm]{img/3/chara.png}
\caption{新規キャラクターの設定画面.}
\label{chara}
\end{center}
\end{figure}

新規キャラクターには, 図\ref{chara}のように, 名前, 役割, 性別, 年齢, 身分・職業, 特徴, 性格という属性を設定することができる. それぞれの項目について『ブラック・ジャック』のキャラクター分析に基づいた選択肢が用意され, その選択肢の中から属性を選ぶことも, 自身で思いついた言葉を当てはめることもできる. ランダム選択機能も備えており, 設定の負担を軽減する.

登場させたい人数分, 新規キャラクターの属性を設定し終えると, 図\ref{chara}に示すように, それぞれのキャラクターに関するプロフィールを生成することができる. ここで言うプロフィールとは, キャラクターの背景や人物像をより詳細に説明する文章のことであり, これは作品上におけるキャラクターの言動に影響する\cite{ref:Mirowski23}. プロフィールは, 設定された属性をもとに生成され, 自身の手で自由に書き換えることもできる. 

以上のような機能により, キャラクター創作を支援する.

\subsubsection{物語展開因子と極性の設定}
物語展開因子とは, 物語の展開に影響を与えるとされる要素のこと\cite{ref:Murai21}である. また, ここで言う極性とは, 物語それぞれの場面における雰囲気を表す.

ここでは, 『ブラック・ジャック』の分析により得られた, 「犯罪と罪の報い」や「夫婦」などの物語展開因子57種の中から, 含めたい要素を複数選択することができる. また, 物語の序盤, 終盤それぞれの極性を, -1(negative), 0(natural), 1(positive)の3つの数字の中から選び設定することができる. 例えば, 序盤に -1, 終盤に 1 を選ぶと, 物語は暗い雰囲気から始まり, 最後には明るい雰囲気で終わるようにすることができる.

\subsubsection{物語構造の設定}
ここでは, 『ブラック・ジャック』の物語構造が, 前のステップで選択された物語展開因子および極性に基づき, 4 から 11 の展開数ごとに複数提案される. ユーザーはその中から生成するプロットの物語構造を一つ設定することができる. これにより, 自身が含めたい要素と適合した, 分析に基づく物語構造をプロットの骨格として定めることができる.
%ここでは, 前のステップで選択された物語展開因子および極性に基づき提案される『ブラック・ジャック』の物語構造の中から, プロットの構造を一つ定めることができる. 
%ここでは, 生成するプロットの物語構造を, 前のステップで選択した物語展開因子および極性に基づき提案される『ブラック・ジャック』の物語構造から一つ選択することができる.

例えば, 提案される物語構造には, 展開数5の場合, 「探索」→「依頼」→「看病・治療」→「情報開示」→「真相発覚」などがある.

\subsubsection{自由指示・時代・場所の設定}
ここでは, 物語の舞台となる時代や場所について設定することができ, また, これまで定めてきた設定を確認することができる. これまでの設定を変更したり, 新たに追加したい要素を自由記述で設定することもできる. これまでの過程で受けた創造性の刺激を自身の言葉でアウトプットする場にもなっている.

\subsubsection{タイトルとあらすじの生成}
ここでは, これまで設定してきた情報をもとに, 1文から2文程度のあらすじとそれに付随するタイトルをそれぞれ10本生成することができる. 

ここで言うあらすじとは, 生成するプロットの元となる短い文章のことを指す.

タイトルのような物語全体に反映するキーワードはクリエーターの創造性を刺激する. また, 1文から2文程度のあらすじは, 物語の方向性に関するイメージを与え, 創造性の支援になると考えられる.

ユーザーは, 生成されたタイトルおよびあらすじから, プロットに活用するものをそのままの文章で選択することも, その内容を自身で加筆修正することもできる. この機能により共創支援を提供する.

\subsubsection{アイデアの生成}
ここでは, これまでの設定をもとに, 具体的な出来事や道具など, プロットの中で用いることができそうな物語全般に関するアイデアを生成することができる. 生成されたアイデアは, これまでの情報をプロットにまで拡張する材料になり得る.

\subsubsection{シーンごとの主観視点等最終設定およびプロットの生成}
ここでは, プロットの各シーンを誰の視点で描くか設定するとともに, これまで定めてきた設定を確認, 修正することができる.
\begin{figure}[htbp]
\begin{center}
\includegraphics[width=40mm]{img/3/settings.png}
\caption{設定の最終確認画面.}
\label{settings}
\end{center}
\end{figure}

シーンごとに主観視点を決められるこの機能は, 本アプリケーションの制作段階で, クリエーターから頂いた意見をもとに実装した機能であり, シーンごとに視点を変えることで, 物語に複雑性を持たせる狙いがある.

ここでの確認, 修正がプロット生成を行う前の最終段階であり, ここで確定させた設定をもとにプロット生成が行われる.

\subsection{プロットの編集}
本アプリケーションでは, プロット生成セクションの実装にとどまらず, LLM とのインタラクティブなやり取りにより, プロットを反復的に洗練させていくことができるプロット編集セクションを実装した. このセクションでは, 既存のプロットを, 新たな指示や要素を加え新しいプロットに書き換えることができ, 新旧プロットのより満足のいく方を採用していくことができる. この過程を反復して行うことでプロットの質を高めることができる. この機能によりクリエーターと LLM とのプロット生成における共創支援を実現する.

%本アプリケーションは, プロットを生成するだけでなく, プロットの改良を目的とした LLM とのインタラクティブなやり取りを可能にするプロット編集セクションを設けている. プロット編集を実行することで, 既存のプロットに, 新たな指示, 要素を加えた, 新たなプロットを生成させることができる. 新しく生成されたプロットと, 既存のプロットを比較し, どちらが良いか選択していくことで, プロットをより良いものにしていくことができる.
実装したプロット編集機能の一覧を以下に示す.
\begin{itemize}
\setlength{\itemsep}{0cm}
\item フリーワードによるやり取り
\item プロット中の単語・文章に関する質問
\item アイデア生成
\item 伏線の処理
\item 一部シーン変更
\item 全体的な雰囲気の変更
\item シーンの入れ替え
\item シーンごとの視点の変更
\item 登場人物の追加
\item 登場人物の役割変更
\item プロットの部分修正
\item 不自然な日本語の修正
\item プロット1シーンの起承転結への分割
\item 手作業での編集
\end{itemize}

\begin{figure}[htbp]
\begin{center}
\includegraphics[width=55mm]{img/3/flow2.png}
\caption{プロット編集セクションにおけるフロー.}
\label{flow2}
\end{center}
\end{figure}

以下では, これらの機能のうち, 実験でよく使われた3つの機能について説明する.

\subsubsection{フリーワードによるやり取り}
この機能を用いることで, ユーザーは生成されたプロットについて自由記述により LLM とやり取りを行うことができる. 自身の言葉でプロットに対する質問を投げかけることや, 新たな指示によりプロットを書き換えることができる.
% \begin{figure}[htbp]
% \begin{center}
% \includegraphics[width=70mm]{img/3/freequestion.png}
% \caption{フリーワードによるやり取り画面.}
% \label{free}
% \end{center}
% \end{figure} 

\subsubsection{プロット中の単語・文章に関する質問}
この機能では, プロットの中における, 説明が不十分であると思われる単語や内容をより深掘りしたいと思う文章について, 詳細な説明と解釈の例を回答として得ることができる. 
\begin{figure}[htbp]
\begin{center}
\includegraphics[width=81mm]{img/3/question.png}
\caption{プロット中の単語・文章に関する質問機能実行の様子. (1): 内容をより詳しく知りたいプロット上の単語や文章をテキストホルダーに入力する. (2): 詳しく聞くを押すと(1)で入力したものを深掘りした文章が生成される. (3): 回答にチェックをつける. (4): プロットに組み込むを押すと(3)でチェックをつけた回答がプロットに反映される.}
\label{question}
\end{center}
\end{figure} 

単に質問を投げかけるだけでは, プロットに書かれていること以上の内容が回答に含まれないため, プロットの範囲外のことも含め回答がなされるよう工夫を施した.

得られた回答を用いて新たなプロットを生成させることもでき, プロットを改善していくことができる.

\subsubsection{一部シーン変更}
この機能は, 選んだシーンの内容を変更することができる機能である. 他のシーンを変えることなく選択したシーンの内容のみを変更することができ, 変更されたシーンが, 矛盾なく前後のシーンと繋がり, プロットとしての一貫性が保たれるよう工夫を施した. 操作は単純で, 変更したいシーンの横にあるチェックボックスにチェックを入れるだけで良い. 操作の簡便性によりクリエーターの負担を軽減する. 

\subsubsection{新旧プロットの比較}
プロット編集セクションでは, 編集機能を用いてプロットの書き換えを行うたび, 新旧のプロットを比較することができる. 新旧プロットは画面上左右に表示され, それらを比較した結果として, 既存のプロットを採用するか, 新しいプロットを採用するか, 2つのプロットを手作業で統合するか選ぶことができる. 繰り返し編集機能を用いることでプロットを満足のいくものへと改良していくことができる.
